\documentclass[UTF8,zihao=false]{ctexart}

\usepackage[a4paper,top=10mm,bottom=10mm,left=12mm,right=12mm]{geometry}
\usepackage{fontspec}
\usepackage{xcolor}
\usepackage{tabularx}
\usepackage{array}
\usepackage{enumitem}
\usepackage{hyperref}

\pagestyle{empty}
\setlength{\parindent}{0pt}
\setlength{\parskip}{0pt}
\setCJKmainfont{Microsoft YaHei}
\setCJKsansfont{Microsoft YaHei}
\setmainfont{Arial}
\setsansfont{Arial}

\definecolor{ink}{HTML}{222222}
\definecolor{muted}{HTML}{5D646C}
\definecolor{accent}{HTML}{1D6F68}
\definecolor{line}{HTML}{D8D8D8}
\hypersetup{colorlinks=true,urlcolor=accent}

\color{ink}
\renewcommand{\arraystretch}{1.08}
\setlist[itemize]{
  leftmargin=1.12em,
  topsep=0pt,
  partopsep=0pt,
  parsep=0pt,
  itemsep=0pt
}

\newcommand{\sectiontitle}[1]{%
  \vspace{3.2pt}
  {\fontsize{10.7pt}{12.4pt}\selectfont\bfseries\color{accent}#1}\par
  \vspace{-3pt}
  \textcolor{line}{\rule{\linewidth}{0.45pt}}
  \vspace{1pt}
}

\newcommand{\entry}[3]{%
  {\fontsize{9.35pt}{10.7pt}\selectfont\bfseries #1}%
  \hfill
  {\fontsize{8.25pt}{9.7pt}\selectfont\color{muted}#2}\par
  {\fontsize{8.35pt}{9.7pt}\selectfont\color{muted}#3}\par
}

\newcommand{\project}[4]{%
  \vspace{1.3pt}
  {\fontsize{8.95pt}{10.4pt}\selectfont\bfseries #1}
  \hfill
  {\fontsize{7.95pt}{9.2pt}\selectfont\color{muted}#2}\par
  {\fontsize{8.05pt}{9.3pt}\selectfont\color{muted}#3}\par
  \begin{itemize}
    #4
  \end{itemize}
}

\newcommand{\skillrow}[2]{%
  {\fontsize{8.35pt}{9.7pt}\selectfont\bfseries #1} &
  {\fontsize{8.35pt}{9.7pt}\selectfont #2}\\
}

\begin{document}

\begin{minipage}[t]{0.44\linewidth}
{\fontsize{23pt}{26pt}\selectfont\bfseries 纪文龙}
\end{minipage}\hfill
\begin{minipage}[t]{0.54\linewidth}
\raggedleft
{\fontsize{8.7pt}{10.4pt}\selectfont
GitHub：\href{https://github.com/JavonLoong}{JavonLoong}\quad
邮箱 / 电话：待补充\\
公开作品：\href{https://javonloong.github.io/RAG/}{RAG 控制台}\quad
\href{https://javonloong.github.io/Mooncake-Modle/}{MoonCake Studio}}
\end{minipage}

\vspace{3pt}
{\fontsize{9.75pt}{11.2pt}\selectfont
清华大学行健书院本科生 \quad | \quad AI 应用工程、工程仿真复刻、产品原型、视觉表达
}

\sectiontitle{个人定位}
{\fontsize{8.45pt}{9.9pt}\selectfont
以 AI 协作完成工程原型和研究交付：能把模糊需求拆成资料检索、数据清洗、Prompt / Skill / SOP、Agent 分工、工具调用、代码实现、测试验证和 PPT / PDF 交付的闭环。
}

\sectiontitle{教育背景}
\entry{清华大学，行健书院，本科生}{2024.09 -- 2028.07}{方向与兴趣：AI 应用工程、知识库 / RAG、工程仿真、交互原型、视觉表达}

\sectiontitle{AI 与工程能力}
\begin{tabularx}{\linewidth}{@{}p{24mm}X@{}}
\skillrow{AI 协作}{Prompt 设计、CoT 任务拆解、Few-shot 风格控制、MCP / 函数调用、Skill / SOP 沉淀、Agent 分工与审核}
\skillrow{工程实现}{Python、FastAPI、ChromaDB、React、Vite、Three.js、HTML/CSS/JavaScript、LaTeX、GitHub Pages}
\skillrow{仿真与数据}{COMSOL 气动极线、AWES 高空风能复刻、数据清洗、向量检索、benchmark、日志与可审计材料组织}
\skillrow{表达交付}{PPT 制作、图像生成与重绘、汇报材料整理、产品概念书、竞品分析、公开网页演示}
\end{tabularx}

\sectiontitle{项目经历}
\project{动力装备知识库控制台}{RAG 与 GraphRAG 预研，公开演示：\href{https://javonloong.github.io/RAG/}{javonloong.github.io/RAG}}
{关键词：知识库构建、向量检索、GraphRAG、后端控制台}
{
  \item 搭建前端上传、入库、检索、benchmark、日志查看控制台；用 FastAPI / Chroma 管理文档切片、索引和运行数据。
  \item 围绕 6098 页动力装备资料完成可审计处理，沉淀 593 条 chunk，推进普通 RAG 到 GraphRAG 与知识图谱构建的前置实验。
}

\project{高空风能 AWES 与 COMSOL 仿真复刻}{工程仿真与软件闭环}
{关键词：工程仿真、COMSOL 建模、数值计算、结果验证、材料组织}
{
  \item 复刻典型飞行轨迹、系留绳张力、功率曲线和 pumping-cycle 能量账本，将 COMSOL 气动极线接入 Python 主仿真。
  \item 搭建 42-case 软件闭环和汇报材料，保留 COMSOL 与 Python 各自承担的证据边界，避免过度包装成完整物理求解。
}

\project{智能电子产品创新实践课程项目}{AI Agent 与 AIGC 设计，复盘视频：\href{https://v.douyin.com/Bx2r1kX1LPo/}{Douyin}}
{关键词：需求拆解、交互流程、Agent 原型、AIGC 3D 设计}
{
  \item 基于科大讯飞星辰平台完成 Agent 创意设计，负责需求拆解、交互流程设计与功能原型搭建，获优秀创意二等奖。
  \item 使用腾讯混元 3D 完成模型生成与设计优化，参与提示词设计、模型效果筛选与展示方案制作，获优秀创意一等奖。
}

\project{美育图像与 PPT 制作}{视觉判断与汇报交付}
{关键词：图像生成、页面重绘、contact sheet、风格一致性、PPT 汇报}
{
  \item 基于“美育 / 薪火计划”等材料完成素材筛选、生成图像、页面重绘、contact sheet 核对和 PPT 整理，重点体现审美判断、版式控制和风格一致性。
}

\project{MoonCake Studio 月饼模具设计器}{3D 参数化建模，公开演示：\href{https://javonloong.github.io/Mooncake-Modle/}{MoonCake Studio}}
{关键词：Three.js、参数化建模、实时预览、STL / GLB 导出}
{
  \item 实现浏览器端 3D 月饼模具自定义工具，支持边缘轮廓、传统花纹、中文刻字、图片浮雕、成品 / 模具双视图预览和 STL / GLB 导出。
}

\end{document}
